\documentclass[11pt]{article}

\usepackage[margin=1in]{geometry}
\usepackage{amsmath,amssymb,amsthm}
\usepackage{mathtools}
\usepackage{listings}
\usepackage{xcolor}
\usepackage{enumitem}
\usepackage{hyperref}
\hypersetup{colorlinks=true,linkcolor=black,urlcolor=blue,citecolor=black}

\newtheorem{theorem}{Theorem}
\newtheorem{lemma}{Lemma}
\newtheorem{proposition}{Proposition}
\theoremstyle{definition}
\newtheorem{definition}{Definition}
\newtheorem{remark}{Remark}

\newcommand{\enc}[1]{[\![ #1 ]\!]}
\newcommand{\wbisim}{\approx}
\newcommand{\Nil}{\mathbf{0}}
\newcommand{\writhe}{\operatorname{w}}
\newcommand{\cost}{\operatorname{c}}

\lstdefinelanguage{rholang}{
  morekeywords={new,in,for,contract,match,if,else,Nil,true,false},
  sensitive=true,
  morecomment=[l]{//},
  morestring=[b]",
}
\lstset{
  language=rholang,
  basicstyle=\ttfamily\small,
  keywordstyle=\color{blue!70!black}\bfseries,
  commentstyle=\color{green!45!black}\itshape,
  stringstyle=\color{red!60!black},
  columns=fullflexible,
  keepspaces=true,
  showstringspaces=false,
  breaklines=true,
  frame=single,
  framesep=4pt,
  xleftmargin=6pt,
}

\title{\bf Knots as Cost--Accounted Rholang Processes:\\
A Dowker--Thistlethwaite--Driven Encoding\\
and Weak--Bisimulation Invariance}
\author{L.~G.~Meredith\\ \small F1R3FLY.io}
\date{\today}

\begin{document}
\maketitle

\begin{abstract}
We revisit a 2005 encoding of knot diagrams as concurrent processes and rebuild
it in core rholang (the reflective higher-order $\rho$-calculus with a built-in
name-restriction operator) equipped with cost accounting. A crossing becomes a
\emph{gate}: the over-strand passes a signal straight through and produces a
private \emph{latch} token; the under-strand is a \emph{join} on its incoming
signal and that latch. Chirality is nothing more than which strand produces the
latch and which joins on it. Driving the construction from a
Dowker--Thistlethwaite (DT) code makes every arc a single-producer /
single-consumer channel, so the arcs \emph{are} the wires and the whole knot is a
flat parallel composition of gates plus injected ``current.'' Under cost
accounting the old \emph{ad hoc} demand for ``enough current'' becomes a solvency
condition on a token ledger, and the surplus current becomes a metered, hence
internal, quantity. We sketch a proof that the encoding, taken as a monoidal
functor from tangles to processes modulo weak bisimulation, is invariant under
the Reidemeister moves, and observe that the cost surplus of the type-I move is
exactly the writhe.
\end{abstract}

\section{Introduction}

The original intuition~\cite{knot2005} was to read a knot diagram as a signalling
network. Each crossing is a small gate with two through-paths, one for the
over-strand and one for the under-strand. A signal on the over-strand passes
straight through; a signal on the under-strand must wait for a \emph{latch} to
open, and the latch is opened by the over-strand's own signal as it passes. Arcs
of the diagram carry signals from one gate port to another. The encoding was
shown to send the Reidemeister moves to weak bisimulations, making
weak-bisimulation class a diagram invariant.

Two ingredients were left implicit. First, because a strand that is over at one
crossing may be under at another, the latch dependencies chain, and a single
circulating signal can deadlock at an under-strand whose latch has not yet been
produced. The original design simply \emph{posited} ``enough current'' to keep the
network live. Second, the encoding was never driven from a machine-readable knot
representation, so it was never mechanically instantiated.

This note repairs both points. We work in core rholang as implemented by the
\texttt{f1r3node-rust} interpreter~\cite{f1r3node}, using the built-in \texttt{new}
(whose several desugarings we do not revisit), the join receipt written with
\texttt{\&}, and a cost semantics that debits a token from a located stack on each
communication. ``Enough current'' becomes a solvency invariant of that ledger,
and the leftover current becomes internal traffic that weak bisimulation is
entitled to erase. We generate the process directly from a DT code, and we recast
the invariance result functorially over tangles so that the equivalence being
preserved is genuinely observable at the tangle boundary.

\section{Core rholang with cost accounting}

We use the fragment of rholang generated by
\[
P,Q \;::=\; \Nil \;\mid\; x!(Q) \;\mid\; \textstyle\sum_i \mathtt{for}(\vec y_i \leftarrow \vec z_i)\{P_i\}
      \;\mid\; P \mid Q \;\mid\; \mathtt{new}\,x\;\mathtt{in}\;\{P\} \;\mid\; {*}x,
\]
with quoting $@P$ and dereferencing ${*}x$ realising reflection. A \emph{join}
receipt $\mathtt{for}(y_1 \leftarrow z_1\ \&\ \cdots\ \&\ y_k \leftarrow z_k)\{P\}$
fires only when a message is simultaneously available on every $z_j$, consuming
one from each; the general receipt is a \texttt{;}-separated list of such
$\&$-joined groups. We rely only on single-channel receives and two-channel
joins.

\paragraph{Cost.}
Fix a commutative monoid of costs $(\mathcal{C},+,0)$ with a phlogiston supply.
The cost semantics attaches to each communication rule a debit: a rendezvous on a
channel $n$ consumes one token from the located stack at $n$ (equivalently, one
unit of phlogiston). A term is \emph{solvent} for a given endowment if no maximal
reduction sequence blocks for want of a token. Write $\cost(P)$ for the multiset
of debits incurred along a run of $P$ to a normal form; for the terms below every
run has the same $\cost$ up to reordering, so $\cost$ is well defined.

\section{The crossing gate}

A crossing has four ports: over-in $oi$, over-out $oo$, under-in $ui$, under-out
$uo$, and a private latch $l$. The gate is
\begin{lstlisting}
new l in {
    for (s <- oi) { oo!(*s) | l!(Nil) }      // over: pass through, produce latch
  | for (s <- ui & _ <- l) { uo!(*s) }       // under: JOIN on signal & latch
}
\end{lstlisting}
The over-strand is the unique producer on $l$; the under-strand is the unique
joiner on $l$. Nothing else distinguishes the two strands, so a crossing and its
mirror differ exactly by swapping producer and joiner. Because the join is atomic,
there is no intermediate state in which the under-signal has been consumed while
the forward is still stalled: the gate either commits both halves or does nothing,
which removes the only awkward interleaving of the nested-receive formulation.

\section{From DT codes to processes}

A DT traversal of an $n$-crossing diagram visits $2n$ \emph{passes}, numbered
$1,\dots,2n$; each crossing is visited twice, receiving one odd and one even
label. The DT code lists, for each odd label $1,3,\dots,2n-1$ in order, the signed
even label sharing that crossing. We use the sign convention: $d_i>0$ iff the
odd-labelled pass is the over-strand (so an alternating diagram has all-positive
code).

Let arc $a_k$ be the segment from pass $k$ to pass $(k \bmod 2n)+1$. Then pass $p$
is \emph{entered} by $a_{p-1}$ (with $a_0 \equiv a_{2n}$) and \emph{exited} by
$a_p$. Two facts make the arcs serve as bare channels:
\begin{proposition}\label{prop:bipartite}
In the encoding of a knot diagram every arc $a_k$ occurs exactly once as the
output port of a gate and exactly once as the input port of a gate.
\end{proposition}
\begin{proof}
Each pass contributes exactly one exit (a send on $a_p$) and one entry (a receive
on $a_{p-1}$). The $2n$ passes are in bijection with the $2n$ arcs on each side,
and a knot diagram is a single closed curve, so the exit/entry incidences are a
perfect matching on arcs.
\end{proof}
Hence there are no forwarder processes: wiring is name-sharing, and each arc is
one metered communication. A short rholang contract \texttt{DTtoKnot} and an even
shorter Python transliteration emit the term; both accompany this note.

\begin{figure}[t]
\begin{lstlisting}
// trefoil 3_1, DT [4,6,2]  (3 crossings, 6 arcs)
new a1, a2, a3, a4, a5, a6 in {
  new l in {   for (s <- a6) { a1!(*s) | l!(Nil) }
             | for (s <- a3 & _ <- l) { a4!(*s) } }   // over a6->a1, under a3->a4
  |
  new l in {   for (s <- a2) { a3!(*s) | l!(Nil) }
             | for (s <- a5 & _ <- l) { a6!(*s) } }   // over a2->a3, under a5->a6
  |
  new l in {   for (s <- a4) { a5!(*s) | l!(Nil) }
             | for (s <- a1 & _ <- l) { a2!(*s) } }   // over a4->a5, under a1->a2
  |
  a1!(Nil) | a2!(Nil) | a3!(Nil) | a4!(Nil) | a5!(Nil) | a6!(Nil)   // current
}
\end{lstlisting}
\caption{The trefoil, generated from its DT code. Each arc \texttt{a$k$} is
produced once and consumed once, per Proposition~\ref{prop:bipartite}.}
\label{fig:trefoil}
\end{figure}

\begin{definition}[Encoding]\label{def:enc}
Given a DT code $d=(d_1,\dots,d_n)$, let $\enc{d}$ be the term
\[
\mathtt{new}\ a_1,\dots,a_{2n}\ \mathtt{in}\ \{\ \big(\textstyle\parallel_{i} G_i\big)\mid I\ \},
\]
where $G_i$ is the gate of crossing $i$ with its ports computed as above and $I =
\parallel_{k=1}^{2n} a_k!(\Nil)$ is the current.
\end{definition}

\section{Current, solvency, and cost}

\begin{lemma}[Solvency]\label{lem:solvent}
For the all-arcs current $I$ of Definition~\ref{def:enc}, $\enc{d}$ is solvent and
deadlock-free: every gate fires exactly once, and a full run consists of $2n$
primary rendezvous ($n$ over-receives and $n$ under-joins).
\end{lemma}
\begin{proof}[Proof sketch]
Each over-receive $\mathtt{for}(s\leftarrow oi)$ has an injected token waiting on
its over-in arc and fires, producing its latch and a forward on its over-out arc.
Once every latch is present, each under-join has its injected under-in token and
its latch, and fires. No receive is ever permanently starved, so there is no
deadlock. The $n$ over-receives and $n$ under-joins are the primary events; the
$2n$ forwarded sends land on arcs whose injected token was already consumed, so
they are residual.
\end{proof}
The residue---one unconsumed token per arc under the all-arcs policy---is exactly
the ``surplus current'' the original design tolerated. Under cost accounting it is
a definite quantity rather than a hope. The \emph{minimal} solvent endowment (the
least current that avoids deadlock) is the genuine optimisation problem hiding
here; it coincides with the phlogiston budget and we leave it open.

\section{Tangles, boundaries, and the encoding functor}

To make ``weak bisimulation'' meaningful we work with \emph{open tangles}. A
tangle $T$ with $p$ inputs and $q$ outputs is a piece of diagram with $p+q$
boundary endpoints; composition glues outputs to inputs, and juxtaposition places
tangles side by side. Tangles form a (braided, framed) monoidal category
$\mathbf{Tang}$ whose morphisms modulo ambient isotopy rel boundary are generated
by cups, caps, the two braidings, and the identity, subject to the Reidemeister
relations~\cite{turaev}.

The encoding of a tangle keeps its boundary arcs \emph{free}: the $p$ input arcs
and $q$ output arcs are not restricted, so they are observable channels. Let
$\mathbf{Proc}_{\wbisim}$ be rholang processes-with-boundary modulo weak
bisimulation $\wbisim$, where $\wbisim$ abstracts the internal (latch and forward)
$\tau$-steps but observes communication on free boundary names. Extend
Definition~\ref{def:enc} to open tangles by supplying local current on internal
arcs only.

\begin{definition}
$\enc{-}\colon \mathbf{Tang}\to \mathbf{Proc}_{\wbisim}$ sends a tangle to its gate
network with free boundary; composition maps to boundary-name identification and
juxtaposition to parallel composition.
\end{definition}

\begin{lemma}[Boundary forwarder]\label{lem:fwd}
For every tangle $T$, $\enc{T}$ is weakly bisimilar to the pure forwarder that
copies each input boundary signal to the output boundary endpoint reached by
following the strand through $T$. All crossing machinery is internal.
\end{lemma}
\begin{proof}[Proof sketch]
By Lemma~\ref{lem:solvent} applied locally, each gate relays its incoming signals
to its outgoing ports through a bounded run of $\tau$-steps, consuming its private
latch. Composing gates along internal arcs composes forwarders; the internal arcs
and latches are restricted, hence their traffic is $\tau$. What remains observable
is precisely the input$\to$output relaying at the free boundary.
\end{proof}

\begin{lemma}[Congruence]\label{lem:cong}
$\wbisim$ is a congruence for parallel composition and name restriction. Hence if
$\enc{T}\wbisim\enc{T'}$ for tangles with equal boundary, then
$\enc{D[T]}\wbisim\enc{D[T']}$ in every diagram context $D[-]$.
\end{lemma}
\begin{proof}
Standard for the $\rho$-calculus: $\mid$ and $\mathtt{new}$ preserve weak
bisimulation. Boundary identification is a substitution of free names, which
$\wbisim$ also respects.
\end{proof}

\section{Weak-bisimulation invariance}

\begin{theorem}[Reidemeister invariance]\label{thm:main}
If diagrams $D$ and $D'$ are related by a finite sequence of Reidemeister moves,
then $\enc{D}\wbisim\enc{D'}$. Consequently $\enc{-}$ descends to a map from knot
types to weak-bisimulation classes.
\end{theorem}
\begin{proof}[Proof sketch]
By Lemma~\ref{lem:cong} it suffices to verify the three moves as local tangle
relations; each relates two tangles with the same boundary.

\emph{Type I (twist).} The kinked $1\!\to\!1$ tangle has a single self-crossing:
one strand is both over and under, its two passes consecutive. The over-pass
produces the latch that the same strand's under-pass immediately joins on, so the
input signal reaches the output through two internal $\tau$-steps. By
Lemma~\ref{lem:fwd} the kink and the bare wire are the same $1\!\to\!1$ forwarder,
hence weakly bisimilar. The cost differs: the kink adds one over-receive and one
under-join, so $\cost(\text{kink})=\cost(\text{wire})+2$.

\emph{Type II (poke).} In the removable configuration one strand $A$ passes over
the other strand $B$ at both crossings. $A$ forwards straight through, producing
two latches; $B$'s signal passes two under-joins, each gated by one of $A$'s
latches. The result is two independent forwarders $A\colon a\!\to\!a'$ and
$B\colon b\!\to\!b'$, identical to the two parallel wires; weak bisimilarity again
follows from Lemma~\ref{lem:fwd}. The cost delta is $+4$ (two over-receives, two
under-joins), local to the two crossings.

\emph{Type III (slide).} Both tangles have three strands and three crossings, and
the move preserves, pairwise, which strand of each crossing is the over-strand, as
well as the input$\to$output connectivity of the three strands. Only the
\emph{location} of the crossings changes. Thus the two encodings induce the same
boundary forwarder and differ only in the internal $\tau$-schedule, so they are
weakly bisimilar. Here $\cost$ is unchanged: both sides have three crossings, i.e.
six primary rendezvous.

Weak bisimilarity of the local tangles lifts to $\enc{D}\wbisim\enc{D'}$ by
Lemma~\ref{lem:cong}, and the Reidemeister moves generate ambient isotopy.
\end{proof}

\begin{remark}[Cost counts the writhe]
Types II and III leave $\cost$ invariant (canceling, resp.\ preserving), while
type I changes it by $+2$ per twist. Hence on regular-isotopy classes the primary
cost is constant, and the only variation across an ambient-isotopy class is
$2\writhe(D)$: \emph{the metered cost surplus of the encoding is twice the
writhe.} If an unframed invariant is wanted, one normalises by this surplus,
exactly as the Kauffman bracket is normalised by a writhe factor to yield the
Jones polynomial. The encoding is thus a regular-isotopy (framed) invariant on the
nose, with the framing read directly off the ledger.
\end{remark}

\section{The Perko pair}\label{sec:perko}

The sharpest classical test for any diagram-level construction is the
\emph{Perko pair}: the two reduced $10$-crossing diagrams Perko~A ($10_{161}$) and
Perko~B (formerly $10_{162}$), listed as distinct knots from Little (1885) through
Rolfsen (1976) until Perko proved them the same knot in
1974~\cite{perko,little}. It is the smallest instance of two reduced diagrams of
the same prime knot that a naive tabulation separated, and---this is the point
that matters for us---the two diagrams have \emph{different writhes}; the pair is
the historical counterexample to Little's 1900 claim that the writhe of a reduced
diagram is a knot invariant~\cite{perkopair}. Perko~A has DT code
\[
[\,4,\;12,\;-16,\;14,\;-18,\;2,\;8,\;-20,\;-10,\;-6\,],
\]
from which $\enc{-}$ produces a $10$-gate, $20$-arc network; every arc is a single
producer/single consumer (Proposition~\ref{prop:bipartite}), and a full run is
$20$ primary rendezvous. The generated term accompanies this note.

The pair probes the encoding from two sides, and it passes both.

\paragraph{Invariance: the theorem must \emph{not} separate them.}
Perko~A and Perko~B are the same knot, hence connected by Reidemeister moves, so
Theorem~\ref{thm:main} gives $\enc{\text{A}}\wbisim\enc{\text{B}}$ as boundary
forwarders, and after closing the boundary
$\enc{\text{A}}_{\mathrm{cl}}\wbisim\enc{\text{B}}_{\mathrm{cl}}$ by the congruence
of $\wbisim$ under restriction (Lemma~\ref{lem:cong}). What makes the Perko pair
famously hard for humans is that Perko's connecting sequence is not monotone in
crossing number: it must pass through diagrams with \emph{more} than ten
crossings. Our proof is indifferent to this. It never tracks crossing number or
any complexity measure along the connecting path; it factors the equivalence
through the three local relations and lifts by congruence. The
crossing-number-increasing detour that defeats a greedy simplifier is invisible
to the argument, so the encoding is not fooled the way the tables were.

\paragraph{Framing: the cost \emph{does} separate them, and should.}
Because the two diagrams have different writhes, the writhe remark gives
\[
\cost(\enc{\text{A}}) - \cost(\enc{\text{B}}) \;=\; 2\bigl(\writhe(\text{A})-\writhe(\text{B})\bigr)\;\neq\;0 .
\]
Thus the encoding assigns the pair the \emph{same} weak-bisimulation class
(correctly: one knot) while recording \emph{distinct} cost ledgers (correctly: two
framings). The unframed invariant is recovered after the writhe normalisation of
the remark; the framed, regular-isotopy reading keeps A and B apart exactly where
they genuinely differ.

\paragraph{Why this is a confirmation, not a threat.}
The metered surplus of the encoding is twice the writhe---precisely the quantity
Little's 1900 ``theorem'' wrongly asserted to be a knot invariant, and precisely
the quantity the Perko pair refutes. Theorem~\ref{thm:main} states the correct
status of that quantity: it is a regular-isotopy (framed) invariant, not an
ambient-isotopy one. The Perko pair instantiates that very distinction at minimal
crossing number. Far from embarrassing the construction, it shows the encoding
cleaves the data at the right joint: knot type lands in the weak-bisimulation
class, framing lands in the cost, and the two are cleanly separated. The atom of
the phenomenon is already the type-I move---a kink on the unknot is the smallest
pair of same-knot diagrams with different writhes, contributing the $+2$ our
ledger charges; the Perko pair is the smallest place where the accumulated effect
distinguishes two reduced diagrams of a nontrivial prime knot.

\section{The Jones polynomial via resolution-contexts}\label{sec:jones}

The gate has a second life. Read dynamically it is a latch; read structurally it is
a four-port \emph{hole}, and the Kauffman bracket is exactly what one computes by
plugging those holes.

\paragraph{Smoothings are contexts.}
A crossing tangle has four boundary ports. There are two planar ways to reconnect
them without a crossing---the two \emph{smoothings}. In process terms each is a
pair of wires, i.e.\ a \emph{context} into which the gate-hole is plugged. Writing
the four ports as top-left, top-right (inputs) and bottom-left, bottom-right
(outputs), abbreviated $\mathrm{TL},\mathrm{TR},\mathrm{BL},\mathrm{BR}$, the two
connectors are
\begin{gather*}
\mathbf{1} = \{\,\mathrm{TL}\text{--}\mathrm{BL},\ \mathrm{TR}\text{--}\mathrm{BR}\,\}
\quad(\text{two through-wires}),\\[2pt]
e = \{\,\mathrm{TL}\text{--}\mathrm{TR},\ \mathrm{BL}\text{--}\mathrm{BR}\,\}
\quad(\text{cup--cap}).
\end{gather*}
Replacing every gate in $\enc{D}$ by one of its two connectors yields a process of
pure wire: a disjoint union of closed loops. A closed loop is a $\tau$-cycle with
no free port; its scalar value is $\delta = -A^2 - A^{-2}$.

\paragraph{The state sum.}
A \emph{state} $s$ assigns a smoothing to each of the $n$ crossings; let $a(s)$ and
$b(s)$ count the $A$- and $B$-smoothings and $|s|$ the number of resulting loops.
The bracket is the sum over context-assignments
\[
\langle D\rangle \;=\; \sum_{s} A^{\,a(s)-b(s)}\;\delta^{\,|s|-1},
\qquad \delta=-A^2-A^{-2},
\]
where the loop count $|s|$ is read directly off the pure-wire process as its number
of $\tau$-cycles. This is precisely a sum over the resolution-contexts of the
gate-holes, weighted by the choice and valued by the closed-loop structure of the
result.

\paragraph{Temperley--Lieb as the algebra of resolution-contexts.}
On a single crossing the two connectors $\mathbf{1}$ and $e$ are the two diagrams of
the Temperley--Lieb algebra $\mathrm{TL}_2$, and Kauffman's skein relation is the
linear expansion
\[
\text{crossing} \;\longmapsto\; A\,\mathbf{1} + A^{-1} e, \qquad e^2 = \delta\, e .
\]
The relation $e^2=\delta e$ is the process fact ``a closed wire loop is worth
$\delta$'': stacking two cup--caps traps one loop. Composition of tangles is
boundary-name identification, juxtaposition is parallel composition, and the closure
of a diagram is the Markov trace that evaluates each remaining loop to $\delta$.
Thus the bracket is the monoidal-linear functor
\[
\langle-\rangle\colon \mathbf{Tang}\longrightarrow \mathbb{Z}[A^{\pm1}]\text{-}\mathbf{Mod},
\qquad \text{factoring through } \mathrm{TL}(\delta),
\]
uniquely determined by its value on the one generating crossing. In the
context-labelled reading, the two smoothings are two context labels; the bracket
sums over context-assignments to the holes, and closed-loop evaluation is the trace.
This is the promised recasting: the Kauffman bracket \emph{is} the linear functor on
the multicategory of resolution-contexts, and $\mathrm{TL}(\delta)$ is the algebra
those contexts generate.

\paragraph{Worked example: the trefoil.}
Present the trefoil as the closure of the $2$-strand braid $\sigma_1^{3}$. Each
$\sigma_1$ maps to $A\,\mathbf{1}+A^{-1}e$ in $\mathrm{TL}_2$; cubing (using
$e^2=\delta e$) gives
\[
(A\,\mathbf{1}+A^{-1}e)^{3} \;=\; A^{3}\,\mathbf{1} \;+\; \bigl(A - A^{-3} + A^{-7}\bigr)\,e .
\]
The trace closure sends $\mathbf{1}\mapsto\delta$ (two loops) and $e\mapsto 1$ (one
loop), so
\[
\langle 3_1\rangle \;=\; A^{3}\delta + \bigl(A - A^{-3}+A^{-7}\bigr)
\;=\; -A^{5} - A^{-3} + A^{-7}.
\]

\paragraph{From bracket to Jones, and the writhe again.}
The bracket is invariant under the type-II and type-III moves but scales by $-A^{\pm3}$
under type-I. The normalisation
\[
f(D) \;=\; (-A^{3})^{-\writhe(D)}\,\langle D\rangle
\]
is invariant under all three moves, and the Jones polynomial is
$V_K(t)=f(D)\big|_{A=t^{-1/4}}$. For the trefoil, $\writhe=+3$ gives
$f=A^{-4}+A^{-12}-A^{-16}$ and hence
\[
V_{3_1}(t) \;=\; -t^{4} + t^{3} + t .
\]
The normalisation exponent is the writhe---the very quantity our cost ledger
records as $2\writhe$ (Section~\ref{sec:perko}). So the bracket-to-Jones passage,
which trades the framed (regular-isotopy) bracket for the ambient-isotopy Jones
invariant, is the \emph{same} writhe quotient as the passage from our
regular-isotopy encoding to an ambient one. One writhe, two faces: on the state-sum
side it is the $(-A^3)^{-\writhe}$ factor; on the process side it is the cost
surplus. The two halves of this note---weak bisimulation with cost, and the
Kauffman state sum---meet exactly at the writhe, and the gate's two readings keep
them apart cleanly: the latch reading carries the dynamics and the current, while
the hole reading, having no over/under, computes the bracket with no current at
all.

\section{Outlook}

\paragraph{A spatial--behavioural classifier.}
Theorem~\ref{thm:main} gives invariance; the converse, distinguishing
non-isotopic knots, is carried by the closed encoding rather than the boundary
forwarder. Closing a tangle into a knot restricts all boundary names, so all
behaviour becomes internal and the discriminating content moves into the
\emph{spatial} structure of the term---the arrangement and nesting of gates and
latches. The logic generated for this structure by the OSLF construction is
adequate by a theorem of the generating algorithm~\cite{oslf}: two closed terms
satisfy the same formulae iff they are bisimilar. The remaining question is purely
one of \emph{faithfulness}---whether the induced bisimulation is as fine as knot
type---not of adequacy. We expect the spatial-behavioural logic to be a sharper
practical classifier than Conway's notation, which is a naming scheme tuned to
rational tangles rather than an intrinsic, isotopy-respecting invariant.

\paragraph{Minimal current.}
The least solvent endowment---the true phlogiston budget of a diagram---may itself
be an isotopy invariant, and its behaviour under the Reidemeister moves is the
natural next computation.

\appendix
\section{Rholang source}\label{app:src}
The listing below is the accompanying file \texttt{knot\_encoding.rho} in full: the
hand-expanded trefoil (Part~1), the parametric \texttt{DTtoKnot} builder (Part~2),
and the Perko~A term of Section~\ref{sec:perko} (Part~3). It targets the
\texttt{f1r3node-rust} rholang interpreter; every send and receive carries a payload
(the interpreter rejects $0$-arity), no receive is guarded, and no arithmetic
appears in a match pattern.

\lstinputlisting[language=rholang,basicstyle=\ttfamily\scriptsize]{knot_encoding.rho}

\begin{thebibliography}{9}
\bibitem{knot2005} L.~G.~Meredith. \emph{Knots as processes} (working notes,
  2005).
\bibitem{f1r3node} F1R3FLY.io. \emph{f1r3node-rust}: a pure-Rust F1R3FLY node with
  a rholang interpreter. \url{https://github.com/F1R3FLY-io/f1r3node-rust}.
\bibitem{oslf} Operational Semantics in Logical Form (OSLF): auto-generation of
  adequate Hennessy--Milner logics for graph-structured lambda theories.
\bibitem{turaev} V.~Turaev. \emph{Operator invariants of tangles and $R$-matrices}
  (for the tangle category and its generators/relations).
\bibitem{perko} K.~A.~Perko, Jr. \emph{On the classification of knots.} Proc.\ Amer.\
  Math.\ Soc.\ \textbf{45} (1974), 262--266.
\bibitem{little} C.~N.~Little. \emph{Non-alternating $\pm$ knots.} Trans.\ Roy.\ Soc.\
  Edinburgh \textbf{39} (1900), 771--778.
\bibitem{perkopair} \emph{Perko pair.} Wikipedia; the two reduced diagrams have
  different writhes, refuting Little's writhe-invariance claim (a Tait conjecture).
\bibitem{kauffman} L.~H.~Kauffman. \emph{State models and the Jones polynomial.}
  Topology \textbf{26} (1987), 395--407.
\bibitem{lickorish} W.~B.~R.~Lickorish. \emph{An Introduction to Knot Theory.}
  Graduate Texts in Mathematics \textbf{175}, Springer, 1997.
\end{thebibliography}

\end{document}
